% =========================================================
% 1. DEFINICIÓN DE COLORES
% =========================================================
\definecolor{primaryColor}{HTML}{1A365D}   % Azul profundo institucional
\definecolor{secondaryColor}{HTML}{2B6CB0} % Azul medio
\definecolor{accentColor}{HTML}{D69E2E}    % Oro/Amarillo matemático
\definecolor{bgColor}{HTML}{F7FAFC}        % Fondo ligeramente gris/azul

% ---------------------------------------------------------
% 2. DISEÑO DE LA PÁGINA DE "PARTE"
% ---------------------------------------------------------
\titleformat{\part}[display]
  {\Huge\bfseries\centering\color{primaryColor}}
  {\Large\color{accentColor}\MakeUppercase{\partname} \thepart}
  {20pt}
  {\Huge}

% ---------------------------------------------------------
% 3. DISEÑO DE CAPÍTULOS (\chapter)
% ---------------------------------------------------------
\titleformat{\chapter}[display]
  {\Huge\bfseries\color{primaryColor}} % Formato del título principal
  {\Large\color{accentColor}\MakeUppercase{\chaptername} \Huge\thechapter} % "CAPÍTULO 1"
  {0.5pt} % Espacio vertical entre el número y el título
  {\Huge} % Tamaño de la fuente del título

% ---------------------------------------------------------
% 4. DISEÑO DE SECCIONES (\section)
% ---------------------------------------------------------
% El formato 'hang' es el estándar (número a la izquierda, título a la derecha)
\titleformat{\section}[hang]
  {\Large\bfseries\color{primaryColor}} % Formato general de la sección
  {\color{accentColor}\thesection} % El número (ej. 1.1) en oro
  {0.5em} % Espacio horizontal entre el número y el texto
  {} % Comandos previos al título (vacío en este caso)

% ---------------------------------------------------------
% 5. DISEÑO DE SUBSECCIONES (\subsection)
% ---------------------------------------------------------
\titleformat{\subsection}[hang]
  {\large\bfseries\color{primaryColor}} % Ligeramente más pequeño que la sección
  {\color{accentColor}\thesubsection} % El número (ej. 1.1.1) en oro
  {0.5em}
  {}

  
% Configuración estética para el código de Python
\definecolor{codegreen}{rgb}{0,0.6,0}
\definecolor{codegray}{rgb}{0.5,0.5,0.5}
\definecolor{codepurple}{rgb}{0.58,0,0.82}
\definecolor{backcolour}{rgb}{0.95,0.95,0.92}

\lstset{
    language=Python,
    backgroundcolor=\color{backcolour},   
    commentstyle=\color{codegreen},
    keywordstyle=\color{magenta},
    numberstyle=\tiny\color{codegray},
    stringstyle=\color{codepurple},
    basicstyle=\ttfamily\small,
    breakatwhitespace=false,         
    breaklines=true,                 
    captionpos=b,                    
    keepspaces=true,                 
    showspaces=false,                
    showstringspaces=false,
    showtabs=false,                  
    tabsize=2,
    % --- ESTA ES LA LÍNEA MÁGICA PARA EL ESPAÑOL ---
    literate={á}{{\'a}}1 {é}{{\'e}}1 {í}{{\'i}}1 {ó}{{\'o}}1 {ú}{{\'u}}1
             {Á}{{\'A}}1 {É}{{\'E}}1 {Í}{{\'I}}1 {Ó}{{\'O}}1 {Ú}{{\'U}}1
             {ñ}{{\~n}}1 {Ñ}{{\~N}}1
}